\documentclass[11pt]{article}
\usepackage{amsmath, amssymb, amsthm}
\usepackage{geometry}
\geometry{letterpaper, margin=1in}
\usepackage{enumitem}
\usepackage{hyperref}

\title{Homework 6.2: Homogeneous Linear Equations with Constant Coefficients}
\author{Tashfeen Omran}
\date{September 2026}

\begin{document}

\maketitle
\tableofcontents
\newpage

\section{Part 1: Comprehensive Introduction, Context, and Prerequisites}

\subsection{Core Concepts}
In the study of differential equations, a \textbf{linear homogeneous differential equation with constant coefficients} of order $n$ takes the form:
\[ a_n y^{(n)} + a_{n-1} y^{(n-1)} + \dots + a_1 y' + a_0 y = 0 \]
where $a_n, a_{n-1}, \dots, a_0$ are real constants and $a_n \neq 0$. The term ``homogeneous'' means that the right-hand side of the equation is zero (there is no external forcing function). The term ``linear'' implies that the function $y$ and its derivatives appear only to the first power and are not multiplied together.

Because the equation is linear and homogeneous, it obeys the \textbf{Superposition Principle}: if $y_1(x), y_2(x), \dots, y_n(x)$ are solutions, then any linear combination $C_1 y_1(x) + C_2 y_2(x) + \dots + C_n y_n(x)$ is also a solution. Our goal is to find $n$ linearly independent solutions to form the \textbf{general solution}.

\subsection{Intuition and Derivation}
How do we solve such an equation? We need a function whose derivatives are constant multiples of the function itself, so that when we plug them into the equation, the terms can cancel out to zero. The natural exponential function, $y = e^{rx}$, has exactly this property: $y' = r e^{rx}$, $y'' = r^2 e^{rx}$, and $y^{(n)} = r^n e^{rx}$.

Substituting $y = e^{rx}$ into the differential equation yields:
\[ a_n r^n e^{rx} + a_{n-1} r^{n-1} e^{rx} + \dots + a_1 r e^{rx} + a_0 e^{rx} = 0 \]
Factoring out $e^{rx}$ (which is never zero), we are left with the \textbf{auxiliary (or characteristic) equation}:
\[ a_n r^n + a_{n-1} r^{n-1} + \dots + a_1 r + a_0 = 0 \]
This transforms our calculus problem into an algebra problem: finding the roots of an $n$-th degree polynomial!

\subsection{Historical Context and Motivation}
The techniques for solving these equations were pioneered in the 18th century by Leonhard Euler and later formalized by Augustin-Louis Cauchy. Euler was investigating the dynamics of physical systems, such as vibrating springs and planetary orbits, which are naturally modeled by second-order linear equations (due to Newton's Second Law, $F=ma$, where acceleration is a second derivative). Euler's brilliant insight was to use exponential functions and his famous formula, $e^{ix} = \cos x + i \sin x$, to handle cases where the polynomial roots were complex numbers, thereby unifying exponential and trigonometric functions.

\subsection{Key Formulas}
Let $P(r) = 0$ be the characteristic equation. The general solution depends on the nature of the roots $r$:
\begin{itemize}
    \item \textbf{Distinct Real Roots:} If $r_1, r_2, \dots, r_k$ are distinct real roots, they contribute:
    \[ y(x) = C_1 e^{r_1 x} + C_2 e^{r_2 x} + \dots + C_k e^{r_k x} \]
    \item \textbf{Repeated Real Roots:} If a real root $r$ has multiplicity $m$, it contributes $m$ linearly independent terms:
    \[ y(x) = C_1 e^{rx} + C_2 x e^{rx} + C_3 x^2 e^{rx} + \dots + C_m x^{m-1} e^{rx} \]
    \item \textbf{Complex Conjugate Roots:} If $r = \alpha \pm i\beta$ are complex roots (which always appear in conjugate pairs for real coefficients), they contribute:
    \[ y(x) = e^{\alpha x} (C_1 \cos(\beta x) + C_2 \sin(\beta x)) \]
    \item \textbf{Repeated Complex Roots:} If $\alpha \pm i\beta$ has multiplicity $m$, we multiply by powers of $x$:
    \[ y(x) = e^{\alpha x} \left( (C_1 + C_2 x + \dots) \cos(\beta x) + (D_1 + D_2 x + \dots) \sin(\beta x) \right) \]
\end{itemize}

\subsection{Prerequisites}
To succeed in this topic, you must be proficient in:
\begin{enumerate}
    \item \textbf{Polynomial Algebra:} Factoring polynomials by grouping, synthetic division, and the Rational Root Theorem.
    \item \textbf{Complex Numbers:} Understanding $i = \sqrt{-1}$ and Euler's Formula.
    \item \textbf{Calculus:} The chain rule, product rule, and derivatives of exponential and trigonometric functions.
\end{enumerate}

\newpage
\section{Part 2: Detailed Homework Solutions}
\textit{Note: A strategic selection of 8 problems has been chosen to cover every distinct archetype present in the homework set, ensuring 100\% conceptual coverage without redundant repetition.}

\subsection{Problem 2: Distinct Real Roots}
\textbf{Find a general solution for the differential equation:} $y''' - 3y'' - y' + 3y = 0$

\textbf{Step 1: Write the auxiliary equation.}
Substitute $y = e^{rx}$ to obtain the characteristic polynomial:
\[ r^3 - 3r^2 - r + 3 = 0 \]

\textbf{Step 2: Factor the polynomial.}
We can factor this by grouping. Group the first two terms and the last two terms:
\[ r^2(r - 3) - 1(r - 3) = 0 \]
Factor out the common binomial $(r - 3)$:
\[ (r^2 - 1)(r - 3) = 0 \]
Factor the difference of squares:
\[ (r - 1)(r + 1)(r - 3) = 0 \]

\textbf{Step 3: Identify the roots.}
The roots are $r_1 = 1$, $r_2 = -1$, and $r_3 = 3$. These are distinct real roots.

\textbf{Step 4: Write the general solution.}
\[ y(x) = c_1 e^x + c_2 e^{-x} + c_3 e^{3x} \]

\subsection{Problem 9: Repeated Real Roots}
\textbf{Find a general solution for the differential equation:} $u''' - 9u'' + 27u' - 27u = 0$

\textbf{Step 1: Write the auxiliary equation.}
\[ r^3 - 9r^2 + 27r - 27 = 0 \]

\textbf{Step 2: Factor the polynomial.}
Recognize that this fits the perfect cube binomial expansion pattern: $(a - b)^3 = a^3 - 3a^2b + 3ab^2 - b^3$.
Here, $a = r$ and $b = 3$. Let's verify:
$r^3 - 3(r^2)(3) + 3(r)(3^2) - 3^3 = r^3 - 9r^2 + 27r - 27$.
Thus, the equation factors perfectly as:
\[ (r - 3)^3 = 0 \]

\textbf{Step 3: Identify the roots.}
The root is $r = 3$ with a multiplicity of $m = 3$.

\textbf{Step 4: Write the general solution.}
Because the root is repeated three times, we must multiply by $x$ and $x^2$ to ensure linear independence (assuming $x$ is the independent variable):
\[ u(x) = c_1 e^{3x} + c_2 x e^{3x} + c_3 x^2 e^{3x} \]

\subsection{Problem 13: Repeated Complex Roots}
\textbf{Find a general solution for the differential equation:} $y^{(4)} + 4y'' + 4y = 0$

\textbf{Step 1: Write the auxiliary equation.}
\[ r^4 + 4r^2 + 4 = 0 \]

\textbf{Step 2: Factor the polynomial.}
This is a quadratic in form. Let $u = r^2$. The equation becomes $u^2 + 4u + 4 = 0$, which factors as $(u + 2)^2 = 0$.
Substituting back $u = r^2$:
\[ (r^2 + 2)^2 = 0 \]

\textbf{Step 3: Identify the roots.}
Solving $r^2 + 2 = 0$ gives $r^2 = -2$, so $r = \pm i\sqrt{2}$.
Because of the square on the outside of the parenthesis, these complex roots have a multiplicity of $m = 2$.
Here, $\alpha = 0$ and $\beta = \sqrt{2}$.

\textbf{Step 4: Write the general solution.}
For the first pair of roots, we get $\cos(\sqrt{2}x)$ and $\sin(\sqrt{2}x)$. For the repeated pair, we multiply by $x$:
\[ y(x) = c_1 \cos(\sqrt{2}x) + c_2 \sin(\sqrt{2}x) + c_3 x \cos(\sqrt{2}x) + c_4 x \sin(\sqrt{2}x) \]

\subsection{Problem 15: Operator Form and Mixed Roots}
\textbf{Find a general solution to the given homogeneous equation:} $(D-1)^2(D+3)(D^2+2D+5)^2[y] = 0$

\textbf{Step 1: Translate the operator form to the auxiliary equation.}
The differential operator $D$ corresponds directly to the variable $r$ in the auxiliary equation:
\[ (r - 1)^2 (r + 3) (r^2 + 2r + 5)^2 = 0 \]

\textbf{Step 2: Find the roots for each factor.}
Factor 1: $(r - 1)^2 = 0 \implies r = 1$ (multiplicity 2).
Factor 2: $(r + 3) = 0 \implies r = -3$ (multiplicity 1).
Factor 3: $(r^2 + 2r + 5)^2 = 0$. Use the quadratic formula for $r^2 + 2r + 5 = 0$:
\[ r = \frac{-2 \pm \sqrt{4 - 20}}{2} = \frac{-2 \pm \sqrt{-16}}{2} = \frac{-2 \pm 4i}{2} = -1 \pm 2i \]
Because of the squared factor, $r = -1 \pm 2i$ has multiplicity 2. Here, $\alpha = -1, \beta = 2$.

\textbf{Step 3: Write the general solution.}
Combine the solutions for all roots:
\begin{align*}
y(x) &= \underbrace{(c_1 + c_2 x)e^x}_{\text{from } r=1} + \underbrace{c_3 e^{-3x}}_{\text{from } r=-3} \\
&\quad + \underbrace{e^{-x}(c_4 \cos 2x + c_5 \sin 2x) + x e^{-x}(c_6 \cos 2x + c_7 \sin 2x)}_{\text{from } r=-1 \pm 2i \text{ (repeated)}}
\end{align*}

\subsection{Problem 19: Initial Value Problem}
\textbf{Solve the initial value problem:} $y''' - 4y'' + 4y' = 0; \quad y(0) = -4, \quad y'(0) = -1, \quad y''(0) = -19$

\textbf{Step 1: Find the general solution.}
Auxiliary equation: $r^3 - 4r^2 + 4r = 0$.
Factor out $r$: $r(r^2 - 4r + 4) = 0 \implies r(r - 2)^2 = 0$.
Roots: $r = 0$ (mult 1), $r = 2$ (mult 2).
General solution: 
\[ y(x) = c_1 e^{0x} + c_2 e^{2x} + c_3 x e^{2x} = c_1 + c_2 e^{2x} + c_3 x e^{2x} \]

\textbf{Step 2: Compute the derivatives.}
\[ y'(x) = 2c_2 e^{2x} + c_3(e^{2x} + 2x e^{2x}) = (2c_2 + c_3)e^{2x} + 2c_3 x e^{2x} \]
\[ y''(x) = 2(2c_2 + c_3)e^{2x} + 2c_3(e^{2x} + 2x e^{2x}) = (4c_2 + 4c_3)e^{2x} + 4c_3 x e^{2x} \]

\textbf{Step 3: Apply initial conditions to form a system of equations.}
$y(0) = c_1 + c_2 = -4$
$y'(0) = 2c_2 + c_3 = -1$
$y''(0) = 4c_2 + 4c_3 = -19$

\textbf{Step 4: Solve the system.}
From the second equation, isolate $c_3$:
\[ c_3 = -1 - 2c_2 \]
Substitute this into the third equation:
\[ 4c_2 + 4(-1 - 2c_2) = -19 \]
\[ 4c_2 - 4 - 8c_2 = -19 \]
\[ -4c_2 = -15 \implies c_2 = \frac{15}{4} \]
Now find $c_3$:
\[ c_3 = -1 - 2\left(\frac{15}{4}\right) = -1 - \frac{15}{2} = -\frac{17}{2} \]
Now find $c_1$:
\[ c_1 + \frac{15}{4} = -4 \implies c_1 = -4 - \frac{15}{4} = -\frac{16}{4} - \frac{15}{4} = -\frac{31}{4} \]

\textbf{Step 5: Write the final particular solution.}
\[ y(x) = -\frac{31}{4} + \frac{15}{4} e^{2x} - \frac{17}{2} x e^{2x} \]

\subsection{Problem 22: System of Differential Equations}
\textbf{Find a general solution for the given linear system using the elimination method:}
1) $\frac{d^2x}{dt^2} - x + 5y = 0$
2) $2x + \frac{d^2y}{dt^2} + 2y = 0$

\textbf{Step 1: Isolate $y$ in the first equation.}
\[ 5y = x - x'' \implies y = \frac{1}{5}x - \frac{1}{5}x'' \]
Differentiate twice to find $y''$:
\[ y'' = \frac{1}{5}x'' - \frac{1}{5}x^{(4)} \]

\textbf{Step 2: Substitute $y$ and $y''$ into the second equation.}
\[ 2x + \left(\frac{1}{5}x'' - \frac{1}{5}x^{(4)}\right) + 2\left(\frac{1}{5}x - \frac{1}{5}x''\right) = 0 \]
Multiply the entire equation by 5 to clear denominators:
\[ 10x + x'' - x^{(4)} + 2x - 2x'' = 0 \]
Combine like terms:
\[ -x^{(4)} - x'' + 12x = 0 \implies x^{(4)} + x'' - 12x = 0 \]

\textbf{Step 3: Solve the resulting 4th-order equation for $x(t)$.}
Auxiliary equation: $r^4 + r^2 - 12 = 0$.
Factor as a quadratic in $r^2$:
\[ (r^2 + 4)(r^2 - 3) = 0 \]
Roots: $r^2 = -4 \implies r = \pm 2i$. And $r^2 = 3 \implies r = \pm \sqrt{3}$.
Solution for $x(t)$:
\[ x(t) = c_1 \cos(2t) + c_2 \sin(2t) + c_3 e^{\sqrt{3}t} + c_4 e^{-\sqrt{3}t} \]

\textbf{Step 4: Find $y(t)$ using the relation from Step 1.}
First, compute $x''(t)$:
\[ x''(t) = -4c_1 \cos(2t) - 4c_2 \sin(2t) + 3c_3 e^{\sqrt{3}t} + 3c_4 e^{-\sqrt{3}t} \]
Now substitute into $y = \frac{1}{5}(x - x'')$:
\[ y(t) = \frac{1}{5} \left[ (c_1 - (-4c_1))\cos(2t) + (c_2 - (-4c_2))\sin(2t) + (c_3 - 3c_3)e^{\sqrt{3}t} + (c_4 - 3c_4)e^{-\sqrt{3}t} \right] \]
\[ y(t) = \frac{1}{5} \left[ 5c_1 \cos(2t) + 5c_2 \sin(2t) - 2c_3 e^{\sqrt{3}t} - 2c_4 e^{-\sqrt{3}t} \right] \]
\[ y(t) = c_1 \cos(2t) + c_2 \sin(2t) - \frac{2}{5}c_3 e^{\sqrt{3}t} - \frac{2}{5}c_4 e^{-\sqrt{3}t} \]

\subsection{Problem 31(a): Higher-Order Cauchy-Euler Equation}
\textbf{Solve:} $x^3 y''' + x^2 y'' - 2x y' + 2y = 0, \quad x > 0$

\textbf{Step 1: Use the substitution $y = x^r$.}
Compute the derivatives:
$y' = r x^{r-1}$
$y'' = r(r-1) x^{r-2}$
$y''' = r(r-1)(r-2) x^{r-3}$

\textbf{Step 2: Substitute into the differential equation.}
\[ x^3 (r(r-1)(r-2) x^{r-3}) + x^2 (r(r-1) x^{r-2}) - 2x (r x^{r-1}) + 2 x^r = 0 \]
Factor out $x^r$:
\[ x^r [ r(r-1)(r-2) + r(r-1) - 2r + 2 ] = 0 \]
Since $x > 0$, $x^r \neq 0$. We expand the polynomial inside the brackets:
\[ r(r^2 - 3r + 2) + r^2 - r - 2r + 2 = 0 \]
\[ r^3 - 3r^2 + 2r + r^2 - 3r + 2 = 0 \]
\[ r^3 - 2r^2 - r + 2 = 0 \]

\textbf{Step 3: Factor the polynomial.}
Factor by grouping:
\[ r^2(r - 2) - 1(r - 2) = 0 \implies (r^2 - 1)(r - 2) = 0 \implies (r - 1)(r + 1)(r - 2) = 0 \]
Roots: $r = 1, -1, 2$.

\textbf{Step 4: Write the general solution.}
For Cauchy-Euler equations, the basis solutions are $x^r$:
\[ y(x) = c_1 x + c_2 x^{-1} + c_3 x^2 \]

\subsection{Problem 35: Vibrating Beam}
\textbf{Solve the homogeneous equation for transverse vibrations of a beam:} $EI y^{(4)} - ky = 0$

\textbf{Step 1: Set up the auxiliary equation.}
Divide by $EI$: $y^{(4)} - \frac{k}{EI} y = 0$.
Let $\omega^4 = \frac{k}{EI}$ for convenience, where $\omega > 0$. The equation is $y^{(4)} - \omega^4 y = 0$.
Auxiliary equation: $r^4 - \omega^4 = 0$.

\textbf{Step 2: Factor to find the roots.}
\[ (r^2 - \omega^2)(r^2 + \omega^2) = 0 \]
\[ (r - \omega)(r + \omega)(r^2 + \omega^2) = 0 \]
Roots: $r = \omega, r = -\omega, r = i\omega, r = -i\omega$.

\textbf{Step 3: Write the general solution.}
Using standard exponentials and trigonometric functions:
\[ y(x) = c_1 e^{\omega x} + c_2 e^{-\omega x} + c_3 \cos(\omega x) + c_4 \sin(\omega x) \]
\textit{Note: In engineering, it is highly common to write the real exponential parts in terms of hyperbolic functions. An equivalent and equally correct answer is:}
\[ y(x) = C_1 \cosh(\omega x) + C_2 \sinh(\omega x) + C_3 \cos(\omega x) + C_4 \sin(\omega x) \]

\newpage
\section{Part 3: In-Depth Analysis of Problems and Techniques}

\subsection{A) Problem Types and General Approach}
\begin{itemize}
    \item \textbf{Standard Constant Coefficient (Problems 2, 9, 13):} The core archetype. The strategy is always to write the auxiliary polynomial $P(r) = 0$ and find its roots. The nature of the roots (real/complex, distinct/repeated) dictates the form of the solution.
    \item \textbf{Operator Form (Problem 15):} These problems are actually easier because the polynomial is already factored for you. The operator $D$ acts exactly like the variable $r$. The strategy is to read the roots and their multiplicities directly from the exponents of the factors.
    \item \textbf{Initial Value Problems (Problem 19):} The strategy is a two-phase process. Phase 1: Find the general solution. Phase 2: Take as many derivatives as necessary, plug in $x=0$, and solve the resulting $n \times n$ system of linear equations for the constants $c_1, \dots, c_n$.
    \item \textbf{Systems of Equations (Problem 22):} The strategy is algebraic elimination. Treat the differential operators algebraically. Solve one equation for one variable (e.g., $y$) in terms of the other ($x$) and its derivatives. Substitute this into the second equation to get a single higher-order equation in one variable.
    \item \textbf{Cauchy-Euler Equations (Problem 31a):} These have variable coefficients (powers of $x$ matching the order of the derivative). The strategy is to guess $y = x^r$ instead of $y = e^{rx}$. This yields a polynomial in $r$, but the polynomial is \textit{not} simply the coefficients of the ODE; you must account for the falling factorials $r(r-1)(r-2)\dots$ from the power rule.
\end{itemize}

\subsection{B) Key Algebraic and Calculus Manipulations}
\begin{itemize}
    \item \textbf{Factoring by Grouping:} Used in Problem 2 ($r^3 - 3r^2 - r + 3$) and Problem 31a. This is the most critical algebraic trick for 3rd-degree polynomials in this chapter. You group terms into pairs, factor out the greatest common factor from each pair, and hope a common binomial emerges.
    \item \textbf{Quadratic Form Substitution:} Used in Problem 13 ($r^4 + 4r^2 + 4 = 0$). By recognizing that the polynomial is a quadratic in disguise (letting $u = r^2$), a 4th-degree polynomial becomes trivially easy to factor.
    \item \textbf{Product Rule for Repeated Roots:} In Problem 19, taking the derivative of $c_3 x e^{2x}$ requires the product rule: $c_3(1 \cdot e^{2x} + x \cdot 2e^{2x})$. Forgetting this is the \#1 cause of errors in Initial Value Problems.
    \item \textbf{Operator Algebra:} In Problem 22, we manipulated derivatives as if they were algebraic variables (e.g., substituting $y'' = \frac{1}{5}x'' - \frac{1}{5}x^{(4)}$). This relies on the linearity of the differential operator.
\end{itemize}

\newpage
\section{Part 4: "Cheatsheet" and Tips for Success}

\subsection{Summary of Important Formulas}
For the auxiliary equation $P(r) = 0$:
\begin{center}
\renewcommand{\arraystretch}{1.5}
\begin{tabular}{|l|l|}
\hline
\textbf{Root Type} & \textbf{Terms in General Solution} \\
\hline
Distinct Real: $r_1, r_2$ & $c_1 e^{r_1 x} + c_2 e^{r_2 x}$ \\
\hline
Repeated Real: $r$ (mult. $m$) & $c_1 e^{rx} + c_2 x e^{rx} + c_3 x^2 e^{rx} + \dots + c_m x^{m-1} e^{rx}$ \\
\hline
Distinct Complex: $\alpha \pm i\beta$ & $e^{\alpha x}(c_1 \cos \beta x + c_2 \sin \beta x)$ \\
\hline
Repeated Complex: $\alpha \pm i\beta$ (mult. $m$) & $e^{\alpha x}(c_1 \cos \beta x + c_2 \sin \beta x) + x e^{\alpha x}(c_3 \cos \beta x + c_4 \sin \beta x) + \dots$ \\
\hline
\end{tabular}
\end{center}

\subsection{Tricks and Shortcuts}
\begin{itemize}
    \item \textbf{The Rational Root Theorem:} If you cannot factor by grouping, list all factors of the constant term divided by factors of the leading coefficient. Test these roots using synthetic division. Once you find one root $r=a$, you can divide the polynomial by $(r-a)$ to reduce its degree.
    \item \textbf{Missing Terms:} If an equation is missing the $y$ term (e.g., $y''' - 4y'' + 4y' = 0$), then $r=0$ is guaranteed to be a root. Factor out $r$ immediately.
    \item \textbf{Operator Form Shortcut:} If you see $(D - a)^k y = 0$, do not expand it! The roots are simply $r = a$ with multiplicity $k$.
\end{itemize}

\subsection{Common Pitfalls to Avoid}
\begin{itemize}
    \item \textbf{The "Missing $x$" Trap:} When roots are repeated, students often forget to multiply the second term by $x$, the third by $x^2$, etc. Writing $y = c_1 e^{2x} + c_2 e^{2x}$ is mathematically equivalent to $y = C e^{2x}$, which means you have lost a dimension of your solution space.
    \item \textbf{Confusing $\alpha$ and $\beta$:} For complex roots $r = \alpha \pm i\beta$, the real part $\alpha$ goes in the exponential $e^{\alpha x}$, and the imaginary part $\beta$ goes inside the trig functions $\cos(\beta x)$. Do not put $i$ inside the sine or cosine!
    \item \textbf{Cauchy-Euler vs. Constant Coefficients:} Do not use the $e^{rx}$ auxiliary equation for equations with $x^n y^{(n)}$ terms. You must use the $y = x^r$ substitution, which generates a different characteristic polynomial.
\end{itemize}

\newpage
\section{Part 5: Conceptual Synthesis and The "Big Picture"}

\subsection{A) Thematic Connections}
The central theme of this topic is \textbf{Linearity and Vector Spaces}. The set of all solutions to an $n$-th order homogeneous linear differential equation forms an $n$-dimensional vector space. 
When we find distinct roots, we are essentially finding a set of \textbf{basis vectors} for this space. The constants $c_1, c_2, \dots, c_n$ are the coordinates of a specific solution vector within that space. The requirement that we multiply repeated roots by $x$ is fundamentally a mechanism to ensure our basis vectors remain \textbf{linearly independent} (i.e., they span the full $n$-dimensional space).

\subsection{B) Forward and Backward Links}
\textbf{Backward Links:} This topic is the direct, logical extension of solving second-order constant-coefficient equations. The only difference is that the algebra is harder (factoring cubic and quartic polynomials instead of just using the quadratic formula). It also relies heavily on the linear algebra concepts of linear independence and the Wronskian.

\textbf{Forward Links:} Mastering this section is absolutely non-negotiable for the rest of the course. 
1. \textbf{Non-Homogeneous Equations:} In the next section, you will solve equations where the right side is not zero (e.g., $L[y] = \sin x$). To do this using the Method of Undetermined Coefficients or Variation of Parameters, you \textit{must} first find the homogeneous solution $y_c$ using the exact techniques learned here.
2. \textbf{Laplace Transforms:} Later, you will use Laplace transforms to solve these same equations. The roots of the auxiliary equation will reappear as the "poles" of the transfer function in the complex $s$-plane.

\newpage
\section{Part 6: Abstract Real-World Modeling (AI, ML, \& Quant Finance)}

\subsection{A) Concrete Scenarios}
\begin{enumerate}
    \item \textbf{Quantitative Finance (Jump-Diffusion Models):} In advanced options pricing, asset prices are often modeled with jumps to account for market crashes. Under Kou's Double Exponential Jump-Diffusion model, the integro-differential equation for the option price can be reduced, via differential operators, to a 4th-order linear homogeneous differential equation with constant coefficients.
    \item \textbf{Time Series Analysis (ARIMA Models):} In algorithmic trading, Autoregressive (AR) models predict future prices based on past prices. The stability and stationarity of an AR($n$) process are determined by the roots of its characteristic polynomial, which is the exact discrete-time analog of the auxiliary equations studied in this chapter.
    \item \textbf{Machine Learning (Momentum Optimization):} Advanced gradient descent algorithms (like Heavy Ball or Nesterov momentum) can be modeled as continuous-time differential equations. Analyzing the convergence rates of these algorithms involves solving 2nd and 3rd-order linear ODEs to ensure the system is critically damped.
\end{enumerate}

\subsection{B) Model Problem Setup: Momentum Gradient Descent}
Let us rigorously model the continuous-time limit of Momentum Gradient Descent on a quadratic loss surface.
Define the parameter space $\theta \in \mathbb{R}$. Let the loss function be a strictly convex quadratic form, $L(\theta) = \frac{1}{2} H \theta^2$, where $H > 0$ is the Hessian (curvature).

The discrete Heavy Ball update rule is:
\[ \theta_{k+1} = \theta_k - \alpha \nabla L(\theta_k) + \beta (\theta_k - \theta_{k-1}) \]
Taking the continuous-time limit as the step size $\alpha \to 0$, this difference equation converges to the following second-order ordinary differential equation:
\[ \ddot{\theta}(t) + \gamma \dot{\theta}(t) + \eta H \theta(t) = 0 \]
where $\gamma$ represents the friction (damping) coefficient and $\eta$ is the learning rate.

This is a homogeneous linear differential equation with constant coefficients!
The characteristic equation is:
\[ r^2 + \gamma r + \eta H = 0 \]
The roots are:
\[ r = \frac{-\gamma \pm \sqrt{\gamma^2 - 4\eta H}}{2} \]

\textbf{Optimization Analysis:}
To ensure the neural network parameters $\theta(t)$ converge to the minimum ($\theta = 0$) as fast as possible without oscillating (which causes instability), we want \textbf{critical damping}. This occurs when the roots are real and repeated, meaning the discriminant is zero:
\[ \gamma^2 - 4\eta H = 0 \implies \gamma = 2\sqrt{\eta H} \]
By analyzing the roots of this characteristic equation, ML researchers can optimally tune the momentum hyperparameter $\gamma$ based on the curvature $H$ of the loss landscape!

\newpage
\section{Part 7: Common Variations and Untested Concepts}

\subsection{Boundary Value Problems (BVPs)}
\textbf{Concept:} The homework focused entirely on Initial Value Problems (IVPs), where all conditions are given at a single point (e.g., $y(0), y'(0)$). In physics and engineering (like the vibrating beam in Problem 35), we often encounter Boundary Value Problems, where conditions are given at \textit{two different points} (e.g., the ends of the beam). BVPs can have unique solutions, infinite solutions, or no solution at all.

\textbf{Example Problem:} Solve $y'' + 4y = 0$ subject to $y(0) = 0$ and $y(\pi) = 0$.

\textbf{Solution:}
1. Characteristic equation: $r^2 + 4 = 0 \implies r = \pm 2i$.
2. General solution: $y(x) = C_1 \cos(2x) + C_2 \sin(2x)$.
3. Apply first boundary condition: $y(0) = C_1 \cos(0) + C_2 \sin(0) = C_1 = 0$.
   So, $y(x) = C_2 \sin(2x)$.
4. Apply second boundary condition: $y(\pi) = C_2 \sin(2\pi) = C_2(0) = 0$.
   Notice that $0 = 0$ is true regardless of the value of $C_2$.
5. \textbf{Conclusion:} The solution is $y(x) = C_2 \sin(2x)$ for \textit{any} real number $C_2$. There are infinitely many solutions! This represents the fundamental standing waves of a string fixed at both ends.

\section{Part 8: Advanced Diagnostic Testing: "Find the Flaw"}

\subsection{Problem 1}
\textbf{Prompt:} Solve $y'' - 6y' + 9y = 0$.
\textbf{Flawed Solution:}
The characteristic equation is $r^2 - 6r + 9 = 0$.
Factoring yields $(r - 3)^2 = 0$, so the roots are $r = 3, 3$.
The general solution is $y(x) = C_1 e^{3x} + C_2 e^{-3x}$.
\textbf{Find the Flaw:}
The flaw is that the student used $e^{-3x}$ for the second root instead of multiplying the repeated root's exponential by $x$.
\textbf{Correct Solution:}
Because $r=3$ is a repeated root, the correct general solution is $y(x) = C_1 e^{3x} + C_2 x e^{3x}$.

\subsection{Problem 2}
\textbf{Prompt:} Solve $y'' + 25y = 0$.
\textbf{Flawed Solution:}
The characteristic equation is $r^2 + 25 = 0$.
Solving for $r$, we get $r^2 = -25$, so $r = \pm 5$.
The general solution is $y(x) = C_1 e^{5x} + C_2 e^{-5x}$.
\textbf{Find the Flaw:}
The flaw is taking the square root of a negative number and ignoring the imaginary unit $i$; the roots should be $\pm 5i$, not $\pm 5$.
\textbf{Correct Solution:}
The roots are $r = \pm 5i$. Here $\alpha = 0$ and $\beta = 5$. The correct general solution is $y(x) = C_1 \cos(5x) + C_2 \sin(5x)$.

\subsection{Problem 3}
\textbf{Prompt:} Solve $y''' - y'' = 0$.
\textbf{Flawed Solution:}
The characteristic equation is $r^3 - r^2 = 0$.
Factoring out $r^2$, we get $r^2(r - 1) = 0$.
The roots are $r = 0$ (multiplicity 2) and $r = 1$.
The general solution is $y(x) = C_1 e^{0x} + C_2 x e^{0x} + C_3 e^x$.
Simplifying, $y(x) = C_1 + C_2 + C_3 e^x$.
\textbf{Find the Flaw:}
The flaw is in the simplification step: $C_2 x e^{0x}$ simplifies to $C_2 x$, not $C_2$.
\textbf{Correct Solution:}
The correct simplification is $y(x) = C_1 + C_2 x + C_3 e^x$.

\subsection{Problem 4}
\textbf{Prompt:} Solve the Cauchy-Euler equation $x^2 y'' - 4x y' + 6y = 0$.
\textbf{Flawed Solution:}
The characteristic equation is formed by taking the coefficients: $r^2 - 4r + 6 = 0$.
Using the quadratic formula: $r = \frac{4 \pm \sqrt{16 - 24}}{2} = 2 \pm i\sqrt{2}$.
The solution is $y(x) = x^2 (C_1 \cos(\sqrt{2} \ln x) + C_2 \sin(\sqrt{2} \ln x))$.
\textbf{Find the Flaw:}
The flaw is using the constant-coefficient characteristic equation $a r^2 + b r + c = 0$ instead of the correct Cauchy-Euler characteristic equation $a r(r-1) + b r + c = 0$.
\textbf{Correct Solution:}
The correct equation is $r(r-1) - 4r + 6 = 0 \implies r^2 - 5r + 6 = 0$.
Factoring yields $(r - 2)(r - 3) = 0$, so $r = 2, 3$.
The correct solution is $y(x) = C_1 x^2 + C_2 x^3$.

\subsection{Problem 5}
\textbf{Prompt:} Solve the IVP $y'' - y = 0, \quad y(0) = 1, \quad y'(0) = 2$.
\textbf{Flawed Solution:}
Characteristic equation $r^2 - 1 = 0 \implies r = \pm 1$.
General solution: $y(x) = C_1 e^x + C_2 e^{-x}$.
Apply $y(0) = 1 \implies C_1 + C_2 = 1$.
Derivative: $y'(x) = C_1 e^x + C_2 e^{-x}$.
Apply $y'(0) = 2 \implies C_1 + C_2 = 2$.
This system ($C_1 + C_2 = 1$ and $C_1 + C_2 = 2$) has no solution.
\textbf{Find the Flaw:}
The flaw is an incorrect derivative of $e^{-x}$; the chain rule requires multiplying by $-1$, so $y'(x) = C_1 e^x - C_2 e^{-x}$.
\textbf{Correct Solution:}
The correct derivative is $y'(x) = C_1 e^x - C_2 e^{-x}$.
The system is $C_1 + C_2 = 1$ and $C_1 - C_2 = 2$.
Adding the equations gives $2C_1 = 3 \implies C_1 = 3/2$.
Subtracting gives $2C_2 = -1 \implies C_2 = -1/2$.
The correct solution is $y(x) = \frac{3}{2} e^x - \frac{1}{2} e^{-x}$.

\end{document}